\documentclass[11pt,a4paper]{article}

\usepackage[utf8]{inputenc}
\usepackage[T1]{fontenc}
% Deutsche Silbentrennung, falls die Sprachdateien vorhanden sind. Ohne sie
% uebersetzt das Dokument trotzdem -- nur die Trennung ist dann englisch.
\IfFileExists{ngerman.ldf}{\usepackage[ngerman]{babel}}{}
\usepackage{amsmath,amssymb}
\usepackage{booktabs}
\usepackage{longtable}
\usepackage{array}
\usepackage[margin=2.5cm]{geometry}
% microtype nur mit skalierbaren Schriften (sonst: "auto expansion is only
% possible with scalable fonts").
\IfFileExists{microtype.sty}{\usepackage[protrusion=true,expansion=false]{microtype}}{}
\usepackage[colorlinks=true,linkcolor=black,urlcolor=blue]{hyperref}

\newcommand{\code}[1]{\texttt{\small #1}}
\newcommand{\kappazwei}{\kappa_2}

\title{FEMBEM: Befunde zur Konvergenz des gekoppelten BEM--FEM-Lösers\\
\large TEAM-28-Levitationsproblem}
\author{FEMBEM, Arbeitsprotokoll}
\date{}

\begin{document}
\maketitle

\begin{abstract}
\noindent
Ausgangspunkt war die Frage, ob ein blocktriangulärer Vorkonditionierer mit
Schur-Näherung $\tilde S = \mu_0 G_{\mathrm{BEM}}$ das Konvergenzproblem des
gekoppelten Lösers behebt. \textbf{Die Antwort ist nein}, und die Ursache ist
inzwischen belegt statt vermutet: die Konditionszahl des Systems beträgt
gemessen $\kappazwei = 1{,}07\cdot10^{6}$ und ist ein Sockel aus der Struktur
des Problems. Unterwegs wurden acht reale Fehler im Bestandscode gefunden und
behoben, darunter zwei mit Wirkung im gesamten CPU-Pfad; einer davon war die
Ursache der \code{NaN}-Kräfte. Sechs Diagnosekennzahlen erwiesen sich als
miskalibriert oder irreführend. Zwölf Regressionstests halten die Befunde fest.
Der beste Stand ist unverändert \code{ILU(0)} mit blockdiagonalem
Vorkonditionierer und \code{-{}-gmres-restart=200}.
\end{abstract}

\tableofcontents

%======================================================================
\section{Kurzfassung der Befunde}

\begin{enumerate}
\item \textbf{Der Operator ist korrekt.} \code{apply()} ist eine feste lineare
  Abbildung (Wiederholbarkeit exakt $0$, Linearität $2{,}7\cdot10^{-16}$,
  Übereinstimmung mit der spaltenweise abgetasteten Matrix
  $7{,}5\cdot10^{-16}$).

\item \textbf{Das System ist lösbar.} Die dichte Zerlegung erreicht ein
  Residuum von $4{,}28\cdot10^{-15}$. Keine Inkonsistenz, kein fehlender
  Bildraum.

\item \textbf{Die Konditionszahl ist der Grund.} Gemessen (nicht geschätzt)
  $\kappazwei = 1{,}07\cdot10^{6}$ mit $\sigma_{\max}=155$,
  $\sigma_{\min}=1{,}45\cdot10^{-4}$. Für $\kappazwei\approx10^{6}$ braucht
  unvorkonditioniertes GMRES grob $\sqrt{\kappazwei}\approx10^{3}$ Iterationen
  je Dekade; beobachtet wurden mit \code{ILU(0)} rund $2500$ im späten Bereich.
  Nichts daran ist rätselhaft.

\item \textbf{Der Sockel ist unvermeidbar.} Über die Zeitschrittweite ergibt
  sich eine U-Kurve mit Minimum bei $\kappazwei\approx10^{6}$; die Vorgabe
  $\Delta t = 10^{-4}$ liegt praktisch darauf.

\item \textbf{Der Engpass liegt nicht im FEM-Block.} Vier unabhängige
  Verbesserungen des FEM-Block-Vorkonditionierers, jede an ihrer eigenen
  Kennzahl deutlich messbar, haben den äußeren Löser nicht bewegt.

\item \textbf{$\tilde S = \mu_0 G_{\mathrm{BEM}}$ trägt nicht.} Das innere CG
  divergiert; mehr Iterationen verschlechtern das Ergebnis.
\end{enumerate}

%======================================================================
\section{Getestete Vorkonditionierer-Varianten}

Alle Angaben: relatives Residuum nach $\approx990$ Iterationen, Zeitschritt~1,
Originalnetz (18\,690 FEM-Kanten, 16\,544 BEM-Freiheitsgrade), CPU-Modus.

\begin{table}[h]
\centering
\begin{tabular}{llcc}
\toprule
Vorkonditionierer & Bemerkung & Restart 50 & Restart 200 \\
\midrule
\code{ILU(0)} blockdiagonal   & Vorgabe                       & $0{,}040$ & $\mathbf{0{,}0154}$ \\
AMS blockdiagonal             & vor den Reparaturen           & $1{,}0$   & --- \\
AMS blockdiagonal             & nach Ausblenden               & $0{,}062$ & $0{,}025$ \\
AMS blockdiagonal             & + $\alpha$-Poisson            & ---       & $0{,}023$ \\
AMS blockdiagonal             & + zeilenrel.\ Aggregation     & ---       & $0{,}026$ \\
\code{ILU(0)} blocktriangulär & unsymm.\ $G_{\mathrm{BEM}}$   & $0{,}105$ & $0{,}0333$ \\
\code{ILU(0)} blocktriangulär & symm.\ $G_{\mathrm{BEM}}$     & $0{,}183$ & --- \\
AMS blocktriangulär           &                               & $0{,}12$  & --- \\
ohne Vorkonditionierer        & Diagnosepfad                  & ---       & --- \\
\bottomrule
\end{tabular}
\caption{Varianten am Originalnetz. Der blockdiagonale \code{ILU(0)} gewinnt
bei jeder Restart-Länge.}
\end{table}

\paragraph{Innerer Löser des blocktriangulären Zweigs.}
Erreichtes relatives Residuum des inneren CG auf $\tilde S$ (Ziel $0{,}01$):

\begin{center}
\begin{tabular}{lcc}
\toprule
Innere Iterationen & erreichtes rel.\ Residuum & äußerer Löser \\
\midrule
$3$   & ---              & $0{,}134$ \\
$10$  & $1{,}1\dots2{,}2$ & $0{,}105$ \\
$40$  & ---              & $0{,}0952$ \\
$100$ & $6{,}2\dots9{,}2$ & $0{,}416$ \\
\bottomrule
\end{tabular}
\end{center}

\noindent
Das innere CG divergiert, und mehr Iterationen machen es schlimmer. Die
Schur-Näherung ist damit für diese Diskretisierung widerlegt --- trotz exakt
symmetrischem $G_{\mathrm{BEM}}$.

\paragraph{Einfluss der Restart-Länge.}
Die Restart-Länge war als $50$ fest eincodiert und dominiert alles andere.
Allein ihre Erhöhung auf $200$ verbessert den Standardlauf von $0{,}040$ auf
$0{,}0154$ --- Faktor $2{,}6$ ohne jede Änderung am Vorkonditionierer. Alle
früheren Vergleiche standen dadurch unter Vorbehalt.

%======================================================================
\section{Geprüfte Hypothesen}

\small
\begin{longtable}{p{4.8cm}p{3.8cm}p{5.4cm}}
\toprule
\textbf{Hypothese} & \textbf{Messung} & \textbf{Ergebnis} \\
\midrule
\endhead
Kopplungsmatrix $T$ wird bei\newline verzerrten Hexaedern\newline unsymmetrisch
  & $\max|T_{rc}-T_{cr}|$ auf 200 zufällig verzerrten Quads
  & \textbf{Widerlegt.} Exakt $0$. $\det J$ erscheint in beiden Einträgen an
    denselben Quadraturpunkten. \\
\addlinespace
Sauter-Schwab reicht am\newline Doppelschichtkern bei\newline gemeinsamer Kante nicht
  & Ordnungserhöhung $N=4\dots24$
  & \textbf{Widerlegt.} Exponentielle Konvergenz bis Maschinengenauigkeit
    ($3{,}8\cdot10^{-4}$ bei $N=4$, $6{,}6\cdot10^{-14}$ bei $N=20$). \\
\addlinespace
Kopplungsterm $T z_K$ trägt\newline Gradientenanteil
  & $\|G^{\!\top}v\|/\|v\|$ für $r_A$ und $T z_K$
  & \textbf{Widerlegt.} $1{,}11$ gegen $1{,}26$ --- praktisch gleich, während
    die Verstärkung um Faktor $40$ differiert. \\
\addlinespace
Kohomologie-Projektion der rechten Seite entfernt den Antrieb
  & Lauf mit \code{-{}-nocohomology}
  & \textbf{Widerlegt.} $\|(I-P)b\|$ von $595$ auf $588$, Konvergenz
    unverändert. \\
\addlinespace
Das System ist inkonsistent
  & dichte Zerlegung
  & \textbf{Widerlegt.} Residuum $4{,}28\cdot10^{-15}$. \\
\addlinespace
\code{apply()} ist keine feste lineare Abbildung
  & drei Stufen, u.\,a.\ $A_{\text{dicht}}v$ gegen $\code{apply}(v)$
  & \textbf{Widerlegt.} $0$; $2{,}7\cdot10^{-16}$; $7{,}5\cdot10^{-16}$. \\
\addlinespace
Die Stagnation kommt von der Netzgröße
  & Netzstudie über vier Größen
  & \textbf{Widerlegt.} $1791$ Unbekannte stagnieren wie $18\,690$. \\
\addlinespace
$\kappazwei$ kommt aus dem Materialkontrast $\sigma/\Delta t$ zu $\nu$
  & $\Delta t$-Durchlauf über fünf Werte
  & \textbf{Teilweise.} U-Kurve: beide Äste real, Minimum bei
    $\kappazwei\approx10^{6}$. \\
\addlinespace
Die ausgeblendeten Identitätszeilen ziehen den AMG-Faktor nach oben
  & Modellproblem mit/ohne Maske
  & \textbf{Widerlegt.} $0{,}2295$ gegen $0{,}2149$ --- sie helfen minimal. \\
\addlinespace
Der skalare AMS-Hilfsraum leidet an Gitteranisotropie
  & Faktor über Zell-Seitenverhältnis
  & \textbf{Bestätigt.} $0{,}774$ bei Seitenverhältnis $4{,}5$ (Median der
    Platte) gegen $0{,}237$ bei $1$. \\
\bottomrule
\caption{Von zehn geprüften Hypothesen haben sich acht als falsch erwiesen.
Entscheidend war nicht das Raten, sondern dass jede Vermutung eine Messung nach
sich zog, die sie widerlegen konnte.}
\end{longtable}

%======================================================================
\section{Behobene Fehler im Bestandscode}

\small
\begin{longtable}{p{3.2cm}p{5.0cm}p{5.8cm}}
\toprule
\textbf{Ort} & \textbf{Mechanismus} & \textbf{Wirkung} \\
\midrule
\endhead
\code{bf::dgemv}
  & liest $y$ auch bei $\beta=0$; die BLAS-Spezifikation verbietet das
  & \textbf{Ursache der \code{NaN}-Kräfte.} $d\_temp\_z$ war uninitialisiert,
    $0\cdot\text{NaN}=\text{NaN}$ machte $16\,544$ von $16\,544$ BEM-Einträgen
    zu \code{NaN}. \\
\addlinespace
\code{bf::dnrm2}
  & skalierte Summation; $\text{scale} < a$ ist für \code{NaN} falsch
  & Ein Vektor aus \emph{lauter} \code{NaN} lieferte die Norm $0{,}0$. GMRES
    meldete \textquotedblleft Converged, Rel.\ Res.\ 0\textquotedblright{} bei \code{NaN}-Kräften. Betraf
    alle Abbruchkriterien im CPU-Pfad. \\
\addlinespace
\code{DeviceArray::allocate}
  & lieferte uninitialisierten Speicher
  & 62 Aufrufstellen, keine mit folgendem \code{fill\_zero()}. Nullt jetzt. \\
\addlinespace
GMRES-Rück\-wärts\-sub\-sti\-tu\-tion
  & \code{y[i] /= H[i][i]} ohne jede Prüfung
  & Bei Stagnation wird $H_{ii}$ winzig, $y_i$ unendlich, die Lösung
    \code{NaN} --- ohne Meldung. \\
\addlinespace
\code{DenseLU}, AMG-Grobgitter
  & Singularitätsprüfung \code{best == 0.0}
  & Ein Pivot von $10^{-300}$ kam durch; ein \code{NaN} ebenso, da
    \code{v > best} für \code{NaN} falsch ist. \\
\addlinespace
\code{evaluate\_quad\_}\\\code{normal}
  & \code{1.0/len} ohne Prüfung
  & Entartetes Quad $\to$ Normale aus \code{Inf}/\code{NaN} $\to$ über die
    ACA-Blöcke in den ganzen Operator. \\
\addlinespace
\code{AcaPlus::compress}
  & Pivotsuche \code{|val| > max}; \code{NaN} wird nie gewählt, aber
    gespeichert
  & Ein \code{NaN} aus dem Kern wäre unbemerkt in $U$/$V$ gewandert. \\
\addlinespace
AMG \code{inv\_diag}, AMS $A_G$
  & \code{(d != 0.0) ? 1/d : 0}; rundungskleine Diagonalen
  & $\|z_A\| = 8{,}6\cdot10^{14}$ bei $\|r_A\|=1$. AMS war unbrauchbar. \\
\addlinespace
\code{build\_auxiliary\_}\\\code{nodal\_matrix}
  & Parameter für rundungskleine Diagonalen mit \code{nullptr} aufgerufen
  & Die aussagekräftige Kennzahl wurde berechnet und weggeworfen. \\
\bottomrule
\caption{Acht Fehler plus eine ungenutzte Diagnose. Sieben davon gehören
derselben Familie an: eine Vergleichsoperation, die \code{NaN} oder eine
rundungskleine Größe nicht sehen kann.}
\end{longtable}

\subsection*{Das wiederkehrende Muster}

Jeder Vergleich mit \code{NaN} ist falsch. Damit gilt:
\[
  \code{if (v > max) max = v;} \qquad\text{sieht \code{NaN} nie,}
\]
und eine \code{min}/\code{max}-Statistik kann \code{NaN} \emph{grundsätzlich}
nicht anzeigen. Ebenso schützt \code{d != 0.0} vor der Division durch die exakte
Null, nicht aber durch $10^{-300}$ --- und bei physikalischen Größen ist die
exakte Null der seltene Fall. Die richtige Prüfung ist stets \emph{relativ zu
einem Maßstab} und schließt Endlichkeit ausdrücklich ein. Ausführlich in
\code{docs/design/nan-and-thresholds.md}.

%======================================================================
\section{Miskalibrierte Diagnosekennzahlen}

Sechs Kennzahlen sahen im Fehlerfall harmlos aus und haben die Suche zeitweise
in die Irre geführt. Alle sind korrigiert.

\begin{center}
\small
\begin{tabular}{p{4.4cm}p{9.4cm}}
\toprule
\textbf{Kennzahl} & \textbf{Problem} \\
\midrule
\code{bf::dnrm2} & Norm $0$ statt \code{NaN} $\Rightarrow$ falsche
  Konvergenzmeldung. \\
\code{min}/\code{max} über \code{D\_scale} & \code{NaN} strukturell unsichtbar.
  \\
Ritz-Schlussfolgerung & $\kappazwei\ge1{,}9\cdot10^{3}$ ist eine \emph{untere}
  Schranke; daraus wurde \textquotedblleft 1000 Iterationen reichen\textquotedblright{} gefolgert.
  Gemessen: $1{,}07\cdot10^{6}$, Faktor $560$ daneben. \\
Operatorkomplexität $<1{,}3$ & Schwelle stammt aus einem Modellproblem mit
  ganz anderer Vergröberungsstrategie. \\
AMS-Warntext & empfahl noch den $\alpha$-Poisson-Operator, der längst Vorgabe
  war. \\
AMG-Konvergenzfaktor & wurde ohne die Maske gemessen, die \code{apply()}
  anwendet. \\
\bottomrule
\end{tabular}
\end{center}

%======================================================================
\section{Messungen}

\subsection{Konditionszahl über die Zeitschrittweite}

Netz \code{team28\_coarse\_12} (4511 gekoppelte Unbekannte), $\sigma_{\max}$
durchweg $155$:

\begin{center}
\begin{tabular}{ccccl}
\toprule
$\Delta t$ [s] & Kontrast $\frac{\sigma/\Delta t}{\nu}$ & $\sigma_{\min}$ &
$\kappazwei$ & \\
\midrule
$10^{-6}$ & $4{,}3\cdot10^{7}$ & $7{,}56\cdot10^{-5}$ & $2{,}05\cdot10^{6}$ & \\
$10^{-5}$ & $4{,}3\cdot10^{6}$ & $1{,}50\cdot10^{-4}$ & $1{,}03\cdot10^{6}$ & Minimum \\
$10^{-4}$ & $4{,}3\cdot10^{5}$ & $1{,}45\cdot10^{-4}$ & $1{,}07\cdot10^{6}$ & Vorgabe \\
$10^{-3}$ & $4{,}3\cdot10^{4}$ & $5{,}89\cdot10^{-5}$ & $2{,}63\cdot10^{6}$ & \\
$10^{-2}$ & $4{,}3\cdot10^{3}$ & $7{,}25\cdot10^{-6}$ & $2{,}14\cdot10^{7}$ & \\
\bottomrule
\end{tabular}
\end{center}

\noindent
Eine U-Kurve, kein Trend. Beide Erklärungen sind richtig, jede in einem Ast:
\begin{itemize}
\item \textbf{Großes $\Delta t$:} der Massenterm $\sigma/\Delta t$ wird klein,
  der Operator nähert sich dem reinen Curl-Curl --- und der ist singulär.
\item \textbf{Kleines $\Delta t$:} der Massenterm wird groß, und der Kontrast
  zwischen Leiter ($\sigma/\Delta t$ riesig) und Luft (null) begrenzt.
\end{itemize}
Am Zeitschritt ist also nichts zu gewinnen.

\subsection{Netzstudie}

Ersatznetze mit polar vernetzten Spulen (damit der eingeprägte azimutale Strom
quellenfrei bleibt) und kartesischer Platte; kein Luftelement, wie im Original.

\begin{center}
\begin{tabular}{crcc}
\toprule
$n$ & Zeilen & $\|G^{\!\top}b\|/\|b\|$ & Residuum nach 2000 It. \\
\midrule
$12$ & $1\,791$  & $2{,}4\cdot10^{-8}$ & $0{,}0113$ \\
$16$ & $3\,447$  & $2{,}7\cdot10^{-8}$ & $0{,}0107$ \\
$24$ & $9\,139$  & $2{,}7\cdot10^{-8}$ & $0{,}0128$ \\
$32$ & $17\,199$ & $3{,}6\cdot10^{-8}$ & $0{,}0194$ \\
\midrule
Original & $18\,690$ & $2{,}1\cdot10^{-4}$ & $\approx0{,}011$ \\
\bottomrule
\end{tabular}
\end{center}

\noindent
Über einen Faktor $10$ in der Problemgröße praktisch netzunabhängig. Ein
Konditionierungsproblem aus der Netzweite müsste mit dem Netz wachsen.

\subsection{AMS-Hilfsräume}

Bezugswert: derselbe AMG erreicht auf dem 3D-Laplace $0{,}30$, netzunabhängig.

\begin{center}
\begin{tabular}{lccc}
\toprule
 & Ausgangszustand & nach Ausblenden & Endstand \\
\midrule
skalarer Hilfsraum      & $0{,}93$ & $0{,}72$ & $\mathbf{0{,}52}$ \\
vektorieller Hilfsraum  & $0{,}93$ & $0{,}93$ & $\mathbf{0{,}62}$ \\
Op.-Komplexität skalar  & $1{,}19$ & $1{,}19$ & $\mathbf{3{,}61}$ \\
Op.-Komplexität vektor. & $1{,}09$ & $1{,}09$ & $\mathbf{3{,}58}$ \\
\midrule
äußerer Löser (Restart 200) & --- & $0{,}025$ & $0{,}026$ \\
\bottomrule
\end{tabular}
\end{center}

\noindent
Die Hilfsräume sind deutlich besser, der äußere Löser bewegt sich nicht. Das
ist der Befund, aus dem folgt, dass der Engpass nicht im FEM-Block liegt.

\subsection{Gitteranisotropie und Stärkemaßstab}

Poisson auf $16\times16\times3$, Faktor je V-Zyklus, Grobgittergröße in
Klammern (fein: $1156$):

\begin{center}
\small
\begin{tabular}{lcccc}
\toprule
Zell-Seitenverh. & $1{,}0$ & $4{,}5$ & $10$ & $30$ \\
\midrule
diagonal, $\theta=0{,}08$    & $0{,}237\,(42)$  & $0{,}774\,(40)$ & $0{,}854\,(40)$ & $0{,}873\,(41)$ \\
diagonal, $\theta=0{,}35$    & $0{,}001\,(1156)$ & $0{,}439\,(289)$ & $0{,}478\,(289)$ & $0{,}777\,(189)$ \\
zeilenrelativ, $\theta=0{,}60$ & $0{,}273\,(21)$ & $\mathbf{0{,}439\,(36)}$ & $0{,}478\,(36)$ & $0{,}490\,(36)$ \\
\bottomrule
\end{tabular}
\end{center}

\noindent
Die Aluminiumplatte hat im Originalnetz ein Zell-Seiten\-ver\-hält\-nis von
$4{,}5$ im Median (Maximum $6{,}1$), die Spulen $1{,}4$ bis $1{,}9$. Der
zeilenrelative Maßstab
\[
  |a_{IJ}| > \theta \max_{K\neq I}|a_{IK}|
\]
ist skaleninvariant und bedient beide Regime mit einer Einstellung, bei
gesunder Vergröberung.

\paragraph{Eine Falle.} Der diagonale Maßstab mit $\theta=0{,}35$ liefert im
isotropen Fall einen Faktor von $0{,}0009$ --- aber das Grobgitter ist so groß
wie das feine. Das ist keine Vergröberung, sondern eine teure direkte Lösung.
Ohne die mitgemessene Grobgittergröße wäre das als Durchbruch durchgegangen.

%======================================================================
\section{Neue Werkzeuge}

\subsection{Kommandozeilenschalter}

Fünf Größen waren fest eincodiert und standen damit jeweils einer Messung im
Weg.

\begin{center}
\small
\begin{tabular}{ll}
\toprule
\code{-{}-gmres-restart=N} & Restart-Länge (war $50$). Dominiert alles. \\
\code{-{}-max-iter=N}      & Iterationsobergrenze (war $1000$). \\
\code{-{}-mesh=DATEI}      & Netzdatei (war \code{team28.vtk}). \\
\code{-{}-dt=X}            & Zeitschrittweite (war $10^{-4}$). \\
\code{-{}-amg-theta=X}     & Stärkeschwelle der Aggregation. \\
\code{-{}-blocktri}        & blocktriangulärer Vorkonditionierer. \\
\code{-{}-pc-inner-iters=N}, \code{-{}-pc-inner-tol=X} & innerer CG-Löser. \\
\code{-{}-direct}          & dichte Zerlegung, Lösbarkeit, $\kappazwei$,
                             Operatorprüfung. \\
\code{-{}-ritz}            & Ritz-Spektrum (war an \code{-{}-verify-field}
                             gekoppelt). \\
\bottomrule
\end{tabular}
\end{center}

\subsection{Regressionstests}

Zwölf Tests, alle ohne GPU lauffähig, alle grün.

\begin{center}
\small
\begin{tabular}{ll}
\toprule
\code{test\_coupling\_sign\_convention} & Operator $(-T)$ und Vorkonditionierer
  $(+T)$ exakt entgegengesetzt \\
\code{test\_block\_triangular\_cons.} & bei $r_K=0$ bitgleich zum
  blockdiagonalen Zweig \\
\code{test\_norm\_nan\_propagation} & \code{dnrm2} verschluckt kein \code{NaN}
  \\
\code{test\_dgemv\_beta\_zero} & \code{dgemv} liest $y$ bei $\beta=0$ nicht \\
\code{test\_singular\_quadrature} & 4-Punkt-Regel trifft die konvergierten
  Werte \\
\code{test\_ams\_scalar\_masking} & Ausblenden begrenzt die Verstärkung \\
\code{test\_gbem\_symmetry} & $\langle x,Gy\rangle = \langle Gx,y\rangle$ \\
\code{test\_amg\_convergence} & AMG löst den 3D-Laplace, netzunabhängig \\
\code{test\_alpha\_poisson} & Vektor-Poisson-Operator als AMG-Grundlage \\
\code{test\_amg\_anisotropy} & Stärkemaßstab auf anisotropen Gittern \\
\code{test\_precond\_end\_to\_end} & Wirkung an der äußeren
  Iterationszahl \\
\code{test\_dense\_direct} & Selbstprüfung von \code{-{}-direct} und der
  $\kappazwei$-Messung \\
\bottomrule
\end{tabular}
\end{center}

\noindent
Drei dieser Tests prüfen ausdrücklich \emph{sich selbst} mit: die
$\kappazwei$-Messung an Matrizen vorgegebener Konditionszahl ($10^{3}$,
$10^{6}$, $10^{9}$, jeweils auf $0{,}0\,\%$ getroffen), der Symmetrietest mit
einer Gegenprobe aus unabhängig gefüllten Blöcken, und der Anisotropietest mit
der Falle des zu großen $\theta$. Ohne solche Gegenproben ist ein auffälliges
Messergebnis nicht von einem Werkzeugfehler zu unterscheiden.

\subsection{Netzgenerator}

\code{tools/make\_coarse\_team28.py} erzeugt grobe Ersatznetze derselben
Geometrie. Zwei Eigenheiten mussten dabei ausdrücklich behandelt werden:

\begin{itemize}
\item \textbf{Keine Luftelemente.} Das Originalnetz hat $4805$ Zellen und
  $1200+1400+2205 = 4805$ --- es enthält ausschließlich die drei Körper, den
  Außenraum macht die BEM.
\item \textbf{Polare Spulen.} Der eingeprägte Spulenstrom ist azimutal; auf
  kartesischem Gitter ist ein Ring eine Treppe, der Strom schließt sich nicht,
  und $\|G^{\!\top}b\|/\|b\|$ steigt von $2{,}1\cdot10^{-4}$ auf $0{,}32$ ---
  das System wird inkonsistent.
\end{itemize}

%======================================================================
\section{Offene Punkte und Empfehlung}

\paragraph{Bester Stand.} \code{ILU(0)} blockdiagonal mit
\code{-{}-gmres-restart=200}: Residuum $0{,}0154$ nach $1000$ Iterationen.

\paragraph{Was noch offen ist.}
\begin{itemize}
\item Der BEM-Block wird von keinem Vorkonditionierer angefasst --- er bekommt
  die Identität. Die blockweise Aufschlüsselung zeigt ihn mit $1{,}8\cdot10^{-4}$
  gegen $1{,}3\cdot10^{-4}$ des FEM-Blocks als den \emph{größeren} Anteil. Ein
  wirksamer Vorkonditionierer für $G_{\mathrm{BEM}}$ (Calderón, hierarchisch)
  wäre der einzige verbliebene Ansatz mit belegter Begründung.
\item Die $3905$ ausgeblendeten Knoten laufen als Identitätszeilen durch die
  gesamte AMG-Hierarchie mit. Das kostet Rechenzeit; auf die Konvergenz wirkt
  es nachweislich nicht.
\item Ein einziges globales $\theta$ kann Platte und Spulen nicht optimal
  zugleich bedienen.
\end{itemize}

\paragraph{Empfehlung.} Vor weiterer Arbeit am Löser sollte die Frage geklärt
werden, \emph{welche Genauigkeit die Kraftwerte tatsächlich brauchen}. Bei
$\kappazwei\approx10^{6}$ ist unterhalb von etwa $10^{-6}$ ohnehin die Grenze
dessen erreicht, was in doppelter Genauigkeit sinnvoll ist. Genügt für eine
belastbare Levitationskraft ein Residuum von $10^{-3}$, ist der aktuelle Stand
bereits ausreichend, und alles Weitere wäre Arbeit an einer Zahl, die niemand
braucht.

%======================================================================
\section*{Methodischer Nachtrag}

Von zehn geprüften Hypothesen erwiesen sich acht als falsch, darunter mehrere,
die zunächst überzeugend aussahen. Zweimal richteten eingebaute
Schutzmaßnahmen mehr Schaden an als der Zustand, den sie abfangen sollten.
Getragen hat nicht das Raten, sondern die Regel, dass jede Vermutung eine
Messung nach sich zog, die sie widerlegen konnte --- und dass jede
Diagnoseausgabe berichtigt wurde, sobald sie einen überholten Stand
weitertrug.

Ein Teil der Arbeit entstand ohne CUDA-Toolkit, geprüft mit \code{g++} gegen
Stub-Header. Das fängt Typfehler und Namensauflösung, aber nicht die strengere
Typherleitung von \code{nvcc} und keinen Gerätecode. Die Grenzen sind in
\code{docs/design/nan-and-thresholds.md} vermerkt.

\end{document}
